% !TEX root=./adl.tex

\section{Introduction}

\begin{frame}{Overall Goals}
  \begin{itemize}
    \item Understand current deep learning state of the art (LLMs, vision, multimodal)
    \item Achieve mastery of deep neural networks
    \item Cover a wide range of topics, fast tempo, in-depth
    \item Coding whenever possible, conceptual and technical depth
    \item After the course you will be well-equipped for a Master's thesis
  \end{itemize}

  \vspace{1em}
  \textbf{Side Effects}
  \begin{itemize}
    \item Earning potential increased (!?)
    \item You will come to appreciate the bitter lesson~\cite{sutton2019bitter}
    \item Ability to follow and discern the scientific literature in deep learning
    \item You will become more productive and iterate faster
    \item You will learn to intuit: what is fundamentally important vs.\ what will fade away?
      Where do performance limitations come from?
    \item Better mental preparation for AGI?
  \end{itemize}
\end{frame}

\begin{frame}{Requirements}
  \begin{columns}
    \begin{column}{0.48\textwidth}
      \textbf{Necessary}
      \begin{itemize}
        \item (Introduction to) Deep Learning
        \item Machine Learning
      \end{itemize}

      \vspace{0.8em}
      \textbf{Desirable}
      \begin{itemize}
          \item Reinforcement Learning (Milica Gasic)
          \item Implementing Transformers Practical (Carol van Niekerk)
      \end{itemize}
    \end{column}
    \begin{column}{0.48\textwidth}
      \textbf{RL Resources}
      \begin{itemize}
        \item \cite{silver2015rl}
        \item \cite{schneider2023rl}
        \item \cite{sutton2018rl}
      \end{itemize}
    \end{column}
  \end{columns}

  \vspace{1em}
  \begin{alertblock}{Interaction}
    Please initiate discussion, pose questions etc.
  \end{alertblock}

  \vspace{1em}
  \begin{alertblock}{Tempo}
    Due to the amount of material covered, I will pursue a higher speed.
  \end{alertblock}

  \vspace{1em}
  \begin{alertblock}{Expectation}
    I expect a big time investment from students. I also invest effort.
  \end{alertblock}

  \vspace{1em}
  \begin{alertblock}{Usage of LLMs}
	  Do not use LLMs for writing code! Looking up function definitions, helping to learn etc.\ is of course alright.
	  Cognitive offloading to LLMs is going to be counterproductive, when you want to learn and understand in depth.
  \end{alertblock}
\end{frame}

\begin{frame}{Organization}
  \begin{columns}
    \begin{column}{0.48\textwidth}
      \textbf{People}
      \begin{itemize}
        \item \textbf{Lecturer:} Paul Swoboda \\
            {\small\email{paul.swoboda@hhu.de}}
        \item \textbf{Assistant:} Pawe\l{} Batorski \\
            {\small\email{pawel.batorski@hhu.de}}
        \item \textbf{Tutor:} Abtin Pourhadi \\
            {\small\email{abtin.pourhadi@hhu.de}}
      \end{itemize}

      \vspace{0.8em}
      \textbf{Administrative questions} \\
      (exam registration, study regulations, etc.)
      \begin{itemize}
        \item Claudia Forstinger \\
          \small\email{claudia.forstinger@hhu.de}
      \end{itemize}
    \end{column}
    \begin{column}{0.48\textwidth}
      \textbf{Communication}
      \begin{itemize}
        \item Primary channel: Discord $\rightarrow$ or would you prefer rocketchat? \\
            %\textcolor{red}{\url{discord.gg/...}}
      \end{itemize}

      \vspace{0.8em}
      \textbf{Materials}
      \begin{itemize}
        \item Github \\
            {\small\textcolor{red}{\url{https://github.com/pawelswoboda/ADL}}}
	\item Please file corrections etc.\ through pull requests
      \end{itemize}

    \end{column}
  \end{columns}
\end{frame}

\begin{frame}{Lecture Modality}
	\begin{itemize}
		\item Assignments will be given each week/second week, depending on scope of exercise. Work on them in groups of three.
		\item Assignments will be corrected, but we will not count points (thank you OpenAI/Anthropic/...)
		\item Two or three mini-exams will determine the grade. One or two exams mid-semester, the final one at the end.
		\item Access to GPUs for completing exercises will be granted on request.
	\end{itemize}
\end{frame}

\begin{frame}{Topics}
  \begin{columns}
    \begin{column}{0.48\textwidth}
      \textbf{Transformers}
      \begin{itemize}
        \item GPT-2 reimplementation ({\`a} la Karpathy)
	\item Recent architectural improvements
	\item Positional Encodings: RoPE, YaRN
	\item Muon optimizer
	\item Training: Pre-training, SFT, RLHF \& post-training
	\item Efficient attention: \{Flash$|$Tind$|$Tree\} attention
	\item KV-Cache: Multi-head latent attention, grouped query attention
	\item Mixture of experts, load balancing
	\item Transformers for computer vision: ViT, DETR, TrackFormer, MaskFormer, SWin, $\ldots$
	\item Multimodal models: CLIP, SigLIP, PaliGemma
      \end{itemize}

      \textbf{Recursive Models}
      \begin{itemize}
	      \item RNNs, LSTM, GRU
	      \item Mamba\{1$|$2\}: Efficient computation with prefix sum \& chunked matrix multiplication
      \end{itemize}
    \end{column}

    \begin{column}{0.48\textwidth}
      \textbf{Training}
      \begin{itemize}
	      \item Efficiency: Mixed precision, kernel fusion 
	      \item Parallelism: Distributed training, Data/model/pipeline parallelism, sharding
          \item Hyperparameters
	      \item Fine-tuning: LoRA, Adapters
	      \item Distillation
	      \item Reinforcement learning: PPO, GRPO
	      \item Scaling laws
      \end{itemize}


      \textbf{Agentic}
      \begin{itemize}
	      \item Prompting: In-context learning, \{chain$|$tree\} of thought
	      \item Memory
	      \item Reasoning
      \end{itemize}

  \begin{alertblock}{Theme}
	  I hope to cover all the essential technical content that allows deep technical understanding of the current iteration of intelligent machines.
  \end{alertblock}
    \end{column}
  \end{columns}
\end{frame}

\begin{frame}{The Bitter Lesson~\cite{sutton2019bitter}}
  \textbf{Core thesis:} General methods leveraging computation beat methods leveraging human knowledge --- and by a large margin.

  \begin{description}[style=nextline]
    \item[Chess] Human-crafted evaluation functions dominated for decades, but massive brute-force search defeated Kasparov in 1997~\cite{campbell2002deep}.
    \item[Go] Initial efforts exploited human knowledge of the game. AlphaGo~Zero learned tabula rasa via self-play + search, surpassing all predecessors~\cite{silver2017mastering}.
    \item[Speech] Hand-engineered phoneme/vocal-tract models lost to statistical HMMs in the 1970s DARPA competition. Deep learning pushed this further~\cite{hinton2012deep}.
    \item[Vision] Edge detectors, SIFT features, generalized cylinders --- all replaced by learned convolutions~\cite{krizhevsky2012imagenet}.
  \end{description}

  \begin{alertblock}{The Lesson}
    \begin{enumerate}
      \item Researchers build domain knowledge into agents --- helps short-term, satisfying personally.
      \item In the long run it plateaus and \emph{inhibits} further progress.
      \item Breakthroughs come from scaling computation via \textbf{search} and \textbf{learning}.
    \end{enumerate}
  \end{alertblock}

  \vspace{0.3em}
  {\small ``We want AI agents that can \emph{discover} like we can, not which contain what we have discovered.'' --- R.~Sutton}
\end{frame}

\begin{frame}{Philosophical/Speculation}
	\begin{itemize}
		\item \underline{AGI}: I think it is a plausible hypothesis that LLMs (or incremental improvements of them) will lead to human-level intelligence in a vast array of settings.
		\item \underline{Timeline:} I think timeslines are on the shorter side.
		\item \underline{Emergence of intelligence}: It is interesting that relatively easy ideas result in emergence of intelligence.
		\item \underline{Ghosts}: LLMs are ghost-like: They are an echo (statistical destillation) of humans organisms, but lack embodiment. Is this a deeper restriction?
		\item \underline{Criticism}: Some people do not believe in LLMs:
			\begin{itemize}
				\item \underline{World models}: Yann LeCun thinks world models are necessary: Causal and physical understanding of the world is not present in LLMs.
				\item \underline{Bitter lesson}: Richard Sutton thinks that LLMs and transformers are not bitter-pilled enough. We only scale through pre- and post-training, in a higly complex manner. A truly bitter-lesson pilled approach would require a single fundamental principle that learns through general interaction with the world, i.e.\ test-time training.
				\item \underline{Jagged Intelligence}: Right now LLMs have jagged intelligence: Superhuman at some complex tasks, struggling with some simple ones. I think that humans suffer from the same, though.
			\end{itemize}
	\end{itemize}
\end{frame}
